\documentclass[11pt,a4paper]{article}
\usepackage[utf8]{inputenc}
\usepackage[margin=1.8cm,headheight=1.5cm,headsep=10pt]{geometry}
\usepackage{graphicx}
\usepackage{fancyhdr}
\usepackage{booktabs}
\usepackage{tabularx}
\usepackage{array}
\usepackage{amsmath,amssymb}
\usepackage{microtype}
\usepackage{hyperref}

\hypersetup{
    colorlinks=true,
    linkcolor=black,
    urlcolor=black,
    citecolor=black
}

\pagestyle{fancy}
\fancyhf{}
\fancyhead[L]{%
  \begin{tabular}{@{}c@{\hspace{7pt}}l@{}}
    \includegraphics[width=1.05cm]{fast_logo.png} & \textbf{FAST NUCES, Islamabad Campus}\\
    & {\small Department of Computer Science -- Fall 2026}
  \end{tabular}%
}
\fancyhead[R]{%
  \begin{tabular}{@{}r@{}}
    \textbf{CS2001: Data Structures}\\
    \textbf{Assignment 2: ImageLab (Phase 1)}\\
    Ahmad Ali \textbar{} 25-i0817 \textbar{} Section E
  \end{tabular}%
}
\fancyfoot[L]{\small FAST NUCES Islamabad}
\fancyfoot[C]{\small CS2001 Data Structures}
\fancyfoot[R]{\small Page \thepage\ of 3}
\renewcommand{\headrulewidth}{0.5pt}
\renewcommand{\footrulewidth}{0.4pt}

\begin{document}

\begin{center}
  {\LARGE \textbf{ImageLab: Phase 1 Technical Report}}\\[4pt]
  {\large \textbf{Custom Dynamic Memory Representations, Core Complexity, and Pipeline Specification}}\\[6pt]
  \textbf{Author:} Ahmad Ali \quad \textbf{Roll Number:} 25-i0817 \quad \textbf{Section:} BS CS-E \quad \textbf{Course:} CS-2001 Data Structures
\end{center}
\vspace{-4pt}
\hrule
\vspace{10pt}

\section{Introduction and Architecture Overview}
The ImageLab system is an interactive image processing studio engineered entirely from first principles without using standard template library (STL) containers, the \texttt{inline} keyword, or C++ \texttt{enum} constructs. Phase 1 establishes the bedrock architecture: raw two-dimensional pointer-based RGB image structures, Malik-compliant templated abstract data types (\texttt{stackType} and \texttt{queueType}), a dual-stack transaction history manager for infinite Undo and Redo, a bounded circular queue execution pipeline, and a dedicated Data Structure Inspector.

\section{Underlying Representations and Justification}

\subsection{Image Memory Representation: 2D Pointer Matrix vs. 1D Flattened Array}
\textbf{Design Decision:} Two-dimensional dynamic array of pixel structs, declared as \texttt{Pixel** imageMatrix}, where each row is independently allocated on the heap as \texttt{new Pixel[width]}.

\begin{itemize}
  \item \textbf{Direct Matrix Correspondence:} Image operations such as 90-degree orthogonal rotations, vertical reflections, and neighborhood kernel convolutions naturally map to Cartesian coordinate pairs $(y, x)$. A 2D pointer structure provides immediate row-level indexing via \texttt{img[y][x]} without requiring repeated multiplications and striding offsets ($\text{index} = y \times \text{width} + x$) inside inner convolution loops.
  \item \textbf{Pointer Swapping in Structural Transformations:} Vertical flipping (\texttt{flipV}) completes in $\mathcal{O}(H)$ time rather than $\mathcal{O}(W \times H)$ because entire horizontal row pointers are swapped (\texttt{std::swap(img[y], img[H - 1 - y])}) without moving or copying individual bytes.
  \item \textbf{Heap Fragmentation Mitigation:} While a contiguous 1D buffer offers marginally better cache spatial locality, a large monolithic allocation frequently fails under memory pressure when heap space is fragmented. Structuring the image as $H$ independent row vectors guarantees reliable dynamic allocation across large canvas buffers.
\end{itemize}

\subsection{History Stacks: Array-Based \texttt{stackType} vs. Linked List Stack}
\textbf{Design Decision:} Bounded array-based dynamic stack (\texttt{stackType<T>}) following Section 5 (Malik interface), using contiguous heap backing and an integer index pointer \texttt{stackTop}.

\begin{itemize}
  \item \textbf{Constant Time Guarantees:} \texttt{push()}, \texttt{pop()}, and \texttt{top()} execute in strict $\mathcal{O}(1)$ worst-case time without heap allocator overhead or pointer dereference stalls.
  \item \textbf{Cache Coherence and Memory Overhead:} An array-based stack avoids 8-byte pointer overhead per node required by linked-list implementations, maximizing L1/L2 cache hit ratios during rapid undo/redo cycles.
\end{itemize}

\subsection{Processing Pipeline: Bounded Circular Array Queue}
\textbf{Design Decision:} Bounded circular array queue with modular arithmetic indexing:
$$\text{qRear} = (\text{qRear} + 1) \pmod{\text{maxQueueSize}}, \quad \text{count} = \text{count} + 1$$

\begin{itemize}
  \item \textbf{Zero Element Shifting:} Unlike naive array queues requiring $\mathcal{O}(N)$ shifting on every dequeue, the circular buffer achieves strict $\mathcal{O}(1)$ enqueue and dequeue operations.
  \item \textbf{Full Queue Handling Policy:} When the circular queue reaches maximum capacity (\texttt{count == maxQueueSize}), new jobs are rejected gracefully with a status code and diagnostic console notification (\texttt{"Pipeline Full: Cannot add job"}). The pipeline does not crash, overwrite unserved tasks, or corrupt memory.
  \item \textbf{Cancellation Without Element Shifting:} Jobs can be cancelled in-place by setting an internal boolean flag \texttt{isCancelled = true}. During execution, the pipeline pops the cancelled job, detects the flag in $\mathcal{O}(1)$, and skips execution without perturbing existing queue indices.
\end{itemize}

\section{Time and Space Complexity Analysis}
Table 1 outlines the formal time and auxiliary space complexities across all core image manipulation algorithms implemented in Phase 1. Let $W$ represent image width, $H$ represent image height, $N = W \times H$ represent total pixel count, and $S$ represent the maximum capacity of history stacks.

\begin{table}[h!]
\centering
\small
\begin{tabularx}{\textwidth}{l c c X}
\toprule
\textbf{Operation} & \textbf{Time Complexity} & \textbf{Auxiliary Space} & \textbf{Engineering Mechanism} \\
\midrule
Load / Save Image & $\mathcal{O}(W \times H)$ & $\mathcal{O}(1)$ & Pure binary/ASCII stream parsing with dynamic 2D row allocation. \\
Crop $(x_1, y_1, x_2, y_2)$ & $\mathcal{O}(W' \times H')$ & $\mathcal{O}(W' \times H')$ & Allocates exact sub-matrix dimensions; copies bounded bounding box. \\
Rotate $90^\circ / 270^\circ$ & $\mathcal{O}(W \times H)$ & $\mathcal{O}(W \times H)$ & Dimension transposition matrix $(H \times W)$; direct coordinate mapping. \\
Rotate $180^\circ$ & $\mathcal{O}(W \times H)$ & $\mathcal{O}(1)$ & In-place double coordinate symmetry swap without extra buffer. \\
Horizontal Flip & $\mathcal{O}(W \times H)$ & $\mathcal{O}(1)$ & In-place swapping of symmetric pixel columns across horizontal rows. \\
Vertical Flip & $\mathcal{O}(H)$ & $\mathcal{O}(1)$ & In-place row pointer swapping without pixel iteration. \\
Brightness Adjustment & $\mathcal{O}(W \times H)$ & $\mathcal{O}(1)$ & Direct linear pixel transform clamped to $[0, 255]$ dynamic range. \\
Grayscale Conversion & $\mathcal{O}(W \times H)$ & $\mathcal{O}(1)$ & NTSC luminance transform: $Y = 0.299R + 0.587G + 0.114B$. \\
Bilinear Resize & $\mathcal{O}(W' \times H')$ & $\mathcal{O}(W' \times H')$ & Bilinear four-point interpolation for aliasing-free down/upscaling. \\
Box Blur ($3 \times 3$) & $\mathcal{O}(W \times H)$ & $\mathcal{O}(W \times H)$ & Frobenius matrix inner product across 8-connected neighborhoods. \\
Bucket Fill (Stack DFS) & $\mathcal{O}(W \times H)$ & $\mathcal{O}(W \times H)$ & Bounded array stack traversal with boolean visited mask. \\
Bucket Fill (Queue BFS) & $\mathcal{O}(W \times H)$ & $\mathcal{O}(W \times H)$ & Bounded circular queue level-order traversal with visited mask. \\
Undo / Redo Dispatch & $\mathcal{O}(W \times H)$ & $\mathcal{O}(W \times H \times S)$ & Full snapshot stack push/pop; deep matrix cloning. \\
Pipeline Enqueue/Dequeue & $\mathcal{O}(1)$ & $\mathcal{O}(1)$ & Modular ring-buffer pointer advancement. \\
\bottomrule
\end{tabularx}
\caption{Algorithmic Time and Space Complexity for Phase 1 Core Operations.}
\end{table}

\section{Empirical Testing and Verification Results}
The implementation has been verified through a dedicated 5-part automated regression testing suite (\texttt{tests/test.cpp}) and interactive graphical inspection:
\begin{enumerate}
  \item \textbf{Stack Primitive Integrity:} Verified LIFO ordering across arbitrary dynamic data types, boundary exception throwing on empty pops, and expansion behavior.
  \item \textbf{Circular Queue Verification:} Verified FIFO order across cyclic wrap-around scenarios, boundary exception handling, and full-queue rejection.
  \item \textbf{Deep Image Operations:} Verified exact mathematical invariants across image transformations:
  $$\text{Rotate}_{90} \circ \text{Rotate}_{90} \circ \text{Rotate}_{90} \circ \text{Rotate}_{90}(I) \equiv I, \quad \text{Flip}_H \circ \text{Flip}_H(I) \equiv I$$
  \item \textbf{Flood Fill Equivalence:} Proved that Stack-based DFS and Queue-based BFS flood fills yield bit-for-bit identical visual output masks while terminating cleanly when source color matches target fill color.
  \item \textbf{Pipeline Concurrency and Cancellation:} Validated bounded job rejection upon reaching maximum capacity (30 items) and verified non-shifting cancellation of jobs.
\end{enumerate}

\section{Summary}
Phase 1 provides an industrial-grade foundation complying with all assignment constraints. All custom Malik template classes, image matrices, and GUI visualizers compile with zero errors and zero warnings, establishing full readiness for Phase 2 memory budgeting and priority scheduling.

\end{document}
